% ===================================================================
%  図表第 1 号 — 学校。— 中学校。— 教室。
%
%  Ceci n'est PAS une transcription : c'est une traduction, faite en
%  2026, en japonais d'aujourd'hui. D'ou trois differences avec
%  texto/io et texto/fr, et trois seulement :
%
%   * pas de \nl ni de \cc — la traduction ne suit aucune ligne
%     imprimee, et les lignes de ce fichier ne sont que du confort ;
%   * pas de \begin{VUpage} — elle n'a ni feuillet ni folio, et la
%     page de lecture ne lui en affiche pas ;
%   * le gras se pose sur le TERME seul, ponctuation dehors.
%
%  Tout le reste est identique, et doit l'etre : les cles « %%K » sont
%  celles de texto/io, macro pour macro pour l'apparat, et les renvois
%  \textsuperscript{(n)} visent les MEMES objets de la planche — ce
%  sont eux qui ouvrent les gros plans.
%
%  L'ido est le texte de base ; le francais sert de controle, et le
%  chinois de calibrage : c'est la seule autre colonne a caracteres
%  pleine chasse, et les deux doivent composer de meme.
%
%  LE GRAS NE SE POSE PAS SUR UN MOT-OUTIL. Le fac-simile ido le fait
%  par accident — « \VUgras{Ica} docisto » — et la page de lecture l'en
%  corrige ; la traduction, elle, n'a pas a repeter l'accident.
%
%  LES NOMS DE PERSONNE SONT NATURALISES, comme dans les neuf autres
%  colonnes, et transcrits en katakana : Ioannes est ジャン, Iakobus
%  est ジャック. Seuls le chien Medor et l'ane Aliboron gardent leur
%  nom.
%
%  LE JAPONAIS MET LE VERBE A LA FIN et postpose ses rapports, comme
%  le hindi : on garde l'ordre des RENVOIS de l'ido, et l'on refait
%  la phrase autour.
%
%  LA PONCTUATION EST PLEINE CHASSE — 。 、 （ ） — comme en chinois.
%  Les chiffres des renvois, eux, restent arabes et demi-chasse : ils
%  numerotent la gravure et doivent se lire d'une colonne a l'autre.
% ===================================================================

%%K t01-apar-0 apar
\VUsaut{2.0mm}
\VUcentre{12.6pt}{120}{説明小冊子}
\VUsaut{3.2mm}
\VUcentre{8.4pt}{160}{附属}
\VUsaut{2.6mm}
\VUtitre{15.6pt}{0}{デルマ補助掛図}
\VUsaut{3.4mm}
\VUornamento{9pt}
\VUsaut{5.0mm}
\VUcentre{13.2pt}{90}{第一部}
\VUsaut{3.0mm}
\VUfilet{16mm}
\VUsaut{5.6mm}

%%K t01-apar-1 apar
\VUtitre{15.0pt}{60}{図表第 1 号}
\VUsaut{1.4mm}
\VUtitre{11.4pt}{0}{学校。--- 中学校。--- 教室。}
\VUsaut{2.2mm}
\VUfilet{20mm}
\VUsaut{6.4mm}

%%K t01-tit-1 sub
\VUpk{10.2pt}{40}{全体の説明。}
\VUsaut{3.0mm}

%%K t01-01-1 p
1. --- この図表が描いているのは、\VUgras{中学校}の\VUgras{教室}である。教室に
は八つの\VUgras{机} \textsuperscript{(18)} と八つの\VUgras{長椅子}
\textsuperscript{(32)} がある。長椅子には一つずつ三人の生徒が坐る。
\VUgras{先生} \textsuperscript{(7)} は\VUgras{椅子} \textsuperscript{(8)} に坐り、自分の
\VUgras{教壇} \textsuperscript{(6)} の上にいる。

%%K t01-02-1 p
2. --- 教壇の左には硝子戸のついた\VUgras{戸棚} \textsuperscript{(15)} が立って
いる。右には\VUgras{黒板} \textsuperscript{(1)} があり、\VUgras{黒板拭き}
\textsuperscript{(2)} と\VUgras{白墨} \textsuperscript{(3)} が添えてある。

%%K t01-02-2 p
教壇の上には\VUgras{掛図} \textsuperscript{(9, 11, 12)} が幾枚かと\VUgras{地図}
\textsuperscript{(10)} が一枚掛かっている。

%%K t01-03-1 p
3. --- 生徒がみな自分の\VUgras{席} \textsuperscript{(17)} にいるわけではない。
ジャン \textsuperscript{(4)} は黒板の前に立っている。黒板拭きを持って、
\VUgras{幾何の図形}を消しているところである。

%%K t01-03-2 p
ピエールは教壇のそばにいる。\VUgras{書き取りの練習} \textsuperscript{(5)} を集め
て、先生のところへ持って行く。

%%K t01-04-1 p
4. --- \VUgras{この}先生は\VUgras{現代語}の先生である。生徒に外国語
\VUgras{（ドイツ語、英語、スペイン語、イタリア語、ロシア語}または
\VUgras{フランス語）}を話し、読み、書くことを教える。

%%K t01-05-1 p
5. --- 教室はたいそう広く、たいそう明るい。奥には二つの\VUgras{欄間}
\textsuperscript{(78)} のついた広い\VUgras{窓}があり、風を入れ光を採るための
\VUgras{硝子戸}がある。

%%K t01-05-2 p
この教室は寒くない。暖めるために、真中に大変よい\VUgras{暖炉}
\textsuperscript{(46)} が置いてあり、その\VUgras{煙突} \textsuperscript{(45)} がついて
いる。

%%K t01-05-3 p
仕切りには\VUgras{衣鉤} \textsuperscript{(58)} が打ちつけてある。そこに生徒の帽子
と外套が掛かっている。

%%K t01-tit-2 sub
\VUsaut{5.2mm}
\VUpk{10.2pt}{40}{運動場。}
\VUsaut{3.6mm}

%%K t01-06-1 p
6. --- \VUgras{時計} \textsuperscript{(63)} は九時五分を指している。

%%K t01-06-2 p
\VUgras{休み時間} \textsuperscript{(76)} はたった今終わり、授業がまた始まる。
休み時間は\VUgras{中庭} \textsuperscript{(75)} で過ごす。何人かの生徒が遊んでいる
のが見える。一人の\VUgras{舎監} \textsuperscript{(74)} が見守っている。

%%K t01-07-1 p
7. --- 中庭にはほかにも何人か見える。\VUgras{視学官} \textsuperscript{(73)}、
\VUgras{校長} \textsuperscript{(71)}、それに\VUgras{教頭} \textsuperscript{(72)} で
ある。

%%K t01-07-2 p
視学官は諸学校を視察し、校長はこの学校を治め、教頭は生徒の学業を監督する。

%%K t01-07-3 p
四人目は\VUgras{用務員} \textsuperscript{(77)} である。欠席者を書きとめる帳面を
脇に抱えている。

%%K t01-08-1 p
8. --- 中庭の奥には\VUgras{体操器械} \textsuperscript{(70)} と\VUgras{手工}
\textsuperscript{(69)} の作業室がある。

%%K t01-08-2 p
図表の右手には\VUgras{音楽}と\VUgras{図画}の\VUgras{教室}がわずかに見える。

%%K t01-noto-1 noto
\VUnotes{143.2mm}{%
(*) \textit{名}については、ラテン語の形でも書き写すことを勧める。数えきれ
ない異形を避けるにはそれしかない。そもそも \textit{Giovanni、Jean、Johann、
Jan、Hans、John、Ivan} のうちどれを選ぶのか。名前のためにわざわざ辞書を一つ
作らねばならぬのか。ラテン語の主格を取って、こう言おう。
\VUgras{Ioannes, Ioanna; Iulius, Iulia; Iakobus; Andreas; Lukas.}
（『完全文法』）。%
}

%%K t01-tit-3 sub
\VUpk{10.2pt}{40}{細かい説明。}
\VUsaut{2.4mm}

%%K t01-tit-4 sub
\VUpk{10.2pt}{40}{よい生徒。}
\VUsaut{3.0mm}

%%K t01-09-1 p
9. --- 教壇の右手、一番目の長椅子の生徒は\VUgras{アンリ}(*)、
\VUgras{エチエンヌ}、\VUgras{シャルル}である。

%%K t01-09-2 p
アンリはたいそう熱心で、たいそう勤勉である。先生の話を聞いている。机の上には
\VUgras{ノート} \textsuperscript{(28)} が一冊、\VUgras{本} \textsuperscript{(29)} が
一冊、\VUgras{辞書} \textsuperscript{(30)} が一冊、\VUgras{筆入れ}
\textsuperscript{(31)} が一つある。彼の本には\VUgras{頁}がたくさんある。どの章にも
\VUgras{段落}が幾つかあり、どの\VUgras{行}にも\VUgras{語}がたくさんある。

%%K t01-09-4 p
隣のエチエンヌ \textsuperscript{(24)} は精を出している。次の時間の課題を自分の
ノートに写しているところである。\VUgras{字}が上手である。彼の本は閉じている。
\VUgras{表紙} \textsuperscript{(27)} しか見えない。そばには彼の\VUgras{コンパス}
\textsuperscript{(26)} が置いてある。

%%K t01-09-5 p
シャルル \textsuperscript{(23)} は手に本を持っている。おさらいのためにそれを開く
ところである。どの生徒とも同じく、彼の前には\VUgras{インク壺}
\textsuperscript{(25)} が一つある。彼は素直である。

%%K t01-10-1 p
10. --- 隣の机には\VUgras{ルイ、アントワーヌ}、\VUgras{ポール}が坐っている。
ルイは自分の\VUgras{課題} \textsuperscript{(22)} を暗誦し、先生の\VUgras{問い}に
答えている。アントワーヌ \textsuperscript{(21)} はたいそう気が散っている。先生は
ときには罰まで与える。ポール \textsuperscript{(20)} はよい仲間で、優しく世話好き
である。たいそう熱心である。

%%K t01-tit-5 sub
\VUsaut{4.2mm}
\VUpk{10.2pt}{40}{わるい生徒。}
\VUsaut{3.2mm}

%%K t01-11-1 p
11. --- アンリの後ろには\VUgras{ジョルジュ} \textsuperscript{(38)} が坐っている。
わるい子ではないが、\VUgras{おしゃべり}で、気が散りやすく、じっとしていられ
ない。彼の\VUgras{軽はずみ}と\VUgras{不従順}が授業中に\VUgras{騒ぎ}を起こし、
きびしい\VUgras{訓戒}を受けて当然のことがしばしばある。彼の\VUgras{下書き帳}
\textsuperscript{(35)} は汚れていて、\VUgras{ペン} \textsuperscript{(34)} で
\VUgras{インクの染み}をつけるものだから、いつも\VUgras{消しゴム}
\textsuperscript{(36)} を使うことになる。彼の\VUgras{鞄} \textsuperscript{(33)} は
破れている。あまりに粗雑だからである。現代語の教室へ\VUgras{ペン軸}
\textsuperscript{(37)} を持って来て何にするのだろう。

%%K t01-11-2 p
隣のエドモン \textsuperscript{(39)} は彼が好きでない。なるほどエドモンは少々怠け者
ではあるが、静かで、じっとしている。むだ話はしない。ただ、聞くかわりに
\VUgras{読本}の\VUgras{お話}を読む方を好むのである。

%%K t01-11-3 p
この長椅子の三人目の生徒は\VUgras{リュドヴィック} \textsuperscript{(40)} である。
精を出している。白い\VUgras{紙}に字を書いている。前には彼の\VUgras{吸取紙}
\textsuperscript{(41)} が広げてある。

%%K t01-12-1 p
12. --- 先生の左手、二番目の長椅子の生徒はたいそう幼い。一人目、小さな
\VUgras{エミール}は、まだ上手に字が書けない。彼の\VUgras{石盤}
\textsuperscript{(42)} と彼の\VUgras{石筆}と彼の\VUgras{海綿} \textsuperscript{(43)}
を持って来ている。

%%K t01-12-2 p
その隣は\VUgras{オーギュスト}と\VUgras{ベルナール} \textsuperscript{(44)} である。
後の者はよく勉強するが、\VUgras{身なり}がたいそうだらしない。

%%K t01-13-1 p
13. --- その後ろには\VUgras{怠け者}が三人いる。\VUgras{アルベール}は課題を覚え
ない。\VUgras{ジャック} \textsuperscript{(47)} は聞くかわりに眠っている。彼の
\VUgras{怠惰}が彼に\VUgras{叱責}を、ときには\VUgras{罰}までも招く。最後に、
\VUgras{クリスチャン} \textsuperscript{(48)} は軽はずみに過ぎる。\VUgras{行い}が
わるく、改めようとも少しもしない。

%%K t01-14-1 p
14. --- ところが隣の生徒たちはよく行いを守っている。\VUgras{エルネスト}
\textsuperscript{(51)} は整頓と規則正しさを好む。いつも自分の\VUgras{定規}
\textsuperscript{(50)} と自分の\VUgras{鉛筆} \textsuperscript{(49)} を使う。彼は
\VUgras{通学生}である。

%%K t01-14-2 p
\VUgras{フランソワ} \textsuperscript{(52)} は\VUgras{寄宿生}である。学校の
\VUgras{制服}を着て、そこで食べ、そこで眠る。よい\VUgras{仲間}である。

%%K t01-14-3 p
モーリスは自分の\VUgras{小刀} \textsuperscript{(56)} で自分の\VUgras{鉛筆}を削って
いる。たいそう念入りである。

%%K t01-14-4 p
本はみな持って来ている。彼の\VUgras{文法書} \textsuperscript{(55)}、彼の
\VUgras{手帳} \textsuperscript{(54)}、それに\VUgras{書き物板} \textsuperscript{(53)}
まである。彼の\VUgras{学校鞄} \textsuperscript{(57)} は後ろの床に置いてある。

%%K t01-14-5 p
教室の奥の二つの机に坐っているのは\VUgras{よい生徒} \textsuperscript{(19)} である。

%%K t01-tit-6 sub
\VUsaut{4.6mm}
\VUpk{10.2pt}{40}{図画と音楽。}
\VUsaut{3.4mm}

%%K t01-15-1 p
15. --- \VUgras{図画教室}では、壁に美しい\VUgras{手本}が掛かり、天井から
\VUgras{笠}のついた\VUgras{ランプ} \textsuperscript{(62)} が下がっている。
\VUgras{先生} \textsuperscript{(60)} はたいそう辛抱強い。生徒の\VUgras{絵}
\textsuperscript{(61)} を直し、\VUgras{胸像} \textsuperscript{(59)} を見ながら写生
することを教える。

%%K t01-16-1 p
16. --- \VUgras{音楽教室}はたいそう面白そうに見える。そこでは\VUgras{楽譜}を
読むこと、\VUgras{五線譜}に書かれた\VUgras{音階} \textsuperscript{(67)} （ド、
レ、ミ、ファ、ソ、ラ、シ\,=\,C. D. E. F. G. A. B.）を歌うこと、
\VUgras{音部記号}を見分けること、美しい\VUgras{曲}を歌うことを習う。楽譜は
\VUgras{譜面台} \textsuperscript{(65)} に置く。生徒の一人が\VUgras{ピアノを弾き}
\textsuperscript{(64)}、\VUgras{音楽の先生} \textsuperscript{(66)} が
\VUgras{拍子をとって} \textsuperscript{(68)} いる。

%%K t01-17-1 p
17. --- \VUgras{手工} \textsuperscript{(69)} の\VUgras{作業室}の隅がわずかに見える。
そこでは子供が木を鉋で削り、鉄を鑢で磨くことを習える。どの中学校にもあるわけ
ではない。

%%K t01-tit-7 sub
\VUsaut{5.0mm}
\VUpk{10.2pt}{40}{理科教室。}
\VUsaut{3.6mm}

%%K t01-18-1 p
18. --- 数学の先生は生徒に\VUgras{幾何、}\VUgras{算術、計算}と
\VUgras{メートル法}を教える。黒板には幾何の主な図形が見える。\VUgras{円}
\textsuperscript{(a)}、\VUgras{正方形} \textsuperscript{(b)}、\VUgras{三角形}
\textsuperscript{(c)}、\VUgras{直線} \textsuperscript{(d)} である。

%%K t01-18-2 p
\VUgras{計算}も一つ見える。それは\VUgras{足し算}でも、\VUgras{引き算}でも、
\VUgras{掛け算}でもない。\VUgras{割り算} \textsuperscript{(e)} である。

%%K t01-19-1 p
19. --- メートル法の掛図には次のものが見分けられる。\VUgras{メートル尺}
\textsuperscript{(a)}、\VUgras{天秤} \textsuperscript{(b)} とその\VUgras{分銅}、
\VUgras{リットル} \textsuperscript{(e)}、\VUgras{デカリットル} \textsuperscript{(f)}、
\VUgras{デシリットル}、\VUgras{立方メートル} \textsuperscript{(d)} である。

%%K t01-19-2 p
生徒は\VUgras{博物} \textsuperscript{(11)} ---動物学 \textsuperscript{(a)}、植物学
\textsuperscript{(b)}、地質学 \textsuperscript{(c)} ---と\VUgras{歴史}
\textsuperscript{(12)} も学ぶ。

%%K t01-19-3 p
戸棚には\VUgras{物理} \textsuperscript{(15)} と\VUgras{化学} \textsuperscript{(16)} の
器械が入っている。

%%K t01-19-4 p
戸棚の上には次のものがある。\VUgras{球} \textsuperscript{(14)}、\VUgras{円錐}、
\VUgras{地球儀} \textsuperscript{(13)} である。

%%K t01-19-5 p
掛図の上には\VUgras{地図}が一枚掛かっていて、世界の五つの部分が描いてある。
\VUgras{ヨーロッパ} \textsuperscript{(g)}、\VUgras{アジア} \textsuperscript{(e)}、
\VUgras{アフリカ} \textsuperscript{(f)}、\VUgras{アメリカ} \textsuperscript{(ab)}、
\VUgras{オセアニア} \textsuperscript{(h)}、それに二つの大洋、\VUgras{太平洋}
\textsuperscript{(c)} と\VUgras{大西洋} \textsuperscript{(d)} である。
\VUsaut{6.0mm}
\VUfilet{22mm}
